\documentclass[a4paper, 11pt]{article}
\usepackage[margin=3cm]{geometry}
\usepackage[]{fontenc}
\usepackage[utf8]{inputenc}
\usepackage[italian]{babel}
\usepackage{geometry}
\geometry{a4paper, top=2cm, bottom=3cm, left=1.5cm, right=1.5cm, heightrounded, bindingoffset=5mm}
\usepackage{amsmath}
\usepackage{amssymb}
\usepackage{gensymb}
\usepackage{graphicx}
\usepackage{psfrag,amsmath,amsfonts,verbatim}
\usepackage{xcolor}
\usepackage{color,soul}
\usepackage{fancyhdr}
\usepackage{indentfirst}
\usepackage{graphicx}
\usepackage{newlfont}
\usepackage{amssymb}
\usepackage{amsmath}
\usepackage{latexsym}
\usepackage{amsthm}
%\usepackage{subfigure}
\usepackage{subcaption}
\usepackage{psfrag}
\usepackage{footnote}
\usepackage{graphics}
\usepackage{color}
\usepackage{hyperref}
\usepackage{tikz}
\usepackage{float}


\usetikzlibrary{snakes}
\usetikzlibrary{positioning}
\usetikzlibrary{shapes,arrows}

	
	\tikzstyle{block} = [draw, fill=white, rectangle, 
	minimum height=3em, minimum width=6em]
	\tikzstyle{sum} = [draw, fill=white, circle, node distance=1cm]
	\tikzstyle{input} = [coordinate]
	\tikzstyle{output} = [coordinate]
	\tikzstyle{pinstyle} = [pin edge={to-,thin,black}]

\newcommand{\courseacronym}{CAT}
\newcommand{\coursename}{Controlli Automatici - T}
\newcommand{\tipology}{C }
\newcommand{\trace}{1}
\newcommand{\projectname}{Controllo di meccanismo non-lineare attuato}
\newcommand{\group}{3 }

%opening
\title{ \vspace{-1in}
		\huge \strut \coursename \strut 
		\\
		\Large  \strut Progetto Tipologia \tipology - Traccia \trace
		\\
		\Large  \strut \projectname\strut
		\\
		\Large  \strut Gruppo \group\strut
		\vspace{-0.4cm}
}
\author{Ester Bravaccini, Camilla Manaresi, Simran Singh}
\date{}

\begin{document}

\maketitle
\vspace{-0.5cm}

Il progetto riguarda il controllo di un meccanismo non-lineare attuato, la cui dinamica viene descritta 

dalla seguente equazione differenziale 
%
\begin{subequations}\label{eq:system}
\begin{align}
    J(\theta)\dot{\omega} = C_m - \beta\omega - k\theta, \\
    \text{dove }
    J(\theta) = J_0 + \sum_{i=1}^{4} J_i \cos(i\theta + \psi_i)
\end{align}
\end{subequations}
%

in cui:
    \begin{itemize}
     \item il momento d'inerzia ridotto al movente \( J(\theta) \) è una funzione periodica di periodo \( 2\pi \);
     \item si assume che \( J(\theta) \) sia già stata ricavata e approssimata a una somma di armoniche;
    \item \( C_m \) rappresenta la coppia generata dal motore elettrico;
    \item si considerano infine anche l'attrito viscoso (coefficiente \( \beta \)) e l'elasticità del disco (coefficiente \( k \)).
\end{itemize}

Si suppone di poter misurare l’angolo \(\theta (t)\) del sistema, ma non la velocità angolare \(\omega (t)\). Inoltre, si 

suppone di poter agire sul sistema tramite la coppia di controllo \( C_m \).

\begin{figure}[H]
    \centering
    \includegraphics[width=0.5\textwidth]{immagini/Figure_0.png}
    \caption{Schema illustrativo del sistema meccanico motorizzato.}
    \label{fig:fig0}
\end{figure} 

I valori numerici dei parametri del sistema sono riportati in Tabella~\ref{tab:system_parameters}.

\begin{table}[H]
\centering
    %\fbox{\parbox{0.35\textwidth}{\vspace{1.9cm} Placeholder \vspace{1.9cm}}}%
\begin{tabular}{|c|c|}
\hline
\textbf{Parametro} & \textbf{Valore} \\
\hline
$k$ & $50$ \\
$\beta$ & $1$ \\
$J_0$ & $1.7$ \\
$J_1$ & $0.1$ \\
$J_2$ & $0.7$ \\
$J_3$ & $0.09$ \\
$J_4$ & $0.02$ \\
$\psi_1$ & $-0.04$ \\
$\psi_2$ & $2.9$ \\
$\psi_3$ & $2.8$ \\
$\psi_4$ & $-2.6$ \\
$\theta_e$ & $\pi/3$ \\
\hline
\end{tabular}
\caption{Parametri del sistema}
\label{tab:system_parameters}
\end{table}

\section{Espressione del sistema in forma di stato e calcolo del sistema linearizzato intorno ad una coppia di equilibrio}

Innanzitutto, esprimiamo il sistema~\eqref{eq:system} nella seguente forma di stato:
%
\begin{subequations}
\begin{align}\label{eq:state_form}
	\dot{x} &= f(x,u)
	\\
	y &= h(x,u).
\end{align}
\end{subequations}
%

Pertanto, procediamo a individuare lo stato $x$, l'ingresso $u$ e l'uscita $y$ del sistema come segue.


Analizziamo che: 
\begin{itemize}
  \item i parametri \(k\), \(\beta\), \(J_0\), \(J_i\) e \(\psi_i\), con \(i = 1,\ldots,4\), sono parametri costanti;
  \item la variabile \(C_m\) è fornita direttamente dal problema come ingresso di controllo e rappresenta la coppia generata dal motore elettrico;
  \item la variabile \(\theta(t)\) rappresenta la posizione angolare del sistema ed è l’unica grandezza misurabile, pertanto costituisce l’uscita del sistema.
\end{itemize}

Sulla base di tali considerazioni, si giunge alle seguenti conclusioni:
\begin{itemize}
    \item la variabile di ingresso \(u(t)\) è rappresentata dalla coppia di controllo \(C_m\), grandezza scalare;
    \item la variabile di stato \(x(t)\) è costituita da due componenti, posizione e velocità angolare, ed è pertanto definita come
    \[
    x =
    \begin{bmatrix}
        x_1 \\
        x_2
    \end{bmatrix},
    \qquad
    \text{con } x_1 := \theta, \; x_2 := \omega;
    \]
    \item la variabile di uscita \(y(t)\) coincide con la grandezza misurata dal sensore ed è pertanto definita come
    \[
    y(t) = x_1 = \theta,
    \]
    che rappresenta l’angolo di rotazione del sistema.
\end{itemize}

Quindi:
%
\begin{align*}
	x := \begin{bmatrix}
		{x}_1
		\\
		{x}_2
	\end{bmatrix} &:= \begin{bmatrix}
	    \theta
        \\
        \omega
	\end{bmatrix}, \quad 
    u :=  C_{m}, \quad 
    y := x_1 := \theta .
\end{align*}
%

Coerentemente con questa scelta, ricaviamo dal sistema~\eqref{eq:system} la seguente espressione per le funzioni $f$ 

ed $h$, in particolare: 
%
\begin{align*}
\dot{x}_1 = f_1(x_1, x_2, u),
\\
\dot{x}_2 = f_2(x_1, x_2, u),
\end{align*}
%
%
\begin{align*}
	f(x,u) &:= \begin{bmatrix}
	   f_1(x_1, x_2,u)
       \\
   f_2(x_1, x_2,u)
	\end{bmatrix} := \begin{bmatrix}
	    \dot{\theta}
        \\
        \dot{\omega} 
	\end{bmatrix} = \begin{bmatrix}
	    x_{2}
        \\
        \frac{u - \beta x_{2} - kx_{1}}{J_0 + \sum_{i=1}^{4} J_i \cos(i x_1 + \psi_i)}
	\end{bmatrix},
	\\
	h(x,u) &:= \theta := x_1.
\end{align*}
%

Una volta calcolate $f$ ed $h$, esprimiamo~\eqref{eq:system} nella seguente forma di stato
%
\begin{subequations}\label{eq:our_system_state_form}
\begin{align}
	\begin{bmatrix}
		\dot{x}_1
		\\
		\dot{x}_2
	\end{bmatrix} &= \begin{bmatrix}
	    x_{2}
        \\
        \frac{u - \beta x_{2} - kx_{1}}{J_0 + \sum_{i=1}^{4} J_i \cos(i x_1 + \psi_i)}
	\end{bmatrix}\label{eq:state_form_1}
	\\
	y &= x_{1}.
\end{align}
\end{subequations}
%

Per trovare la coppia di equilibrio in sistemi LTI continui, se vale $f(x_e, u_e) = 0$, allora $(x_e, u_e)$ è una 

coppia di equilibrio. Il progetto in analisi rientra nel caso riportato.
\\


Considerando il valore di equilibrio dato $\theta_{e} = x_{1e} =\pi/3$, andiamo a risolvere il seguente sistema di 

equazioni: 

%
\begin{align} \begin{bmatrix}
		\dot{x}_1
		\\
		\dot{x}_2
	\end{bmatrix} = 
	\begin{bmatrix}
	    0
        \\
        0
	\end{bmatrix} = \begin{bmatrix}
	    x_{2e}
        \\
        \frac{u_e - \beta x_{2e}-kx_{1e}}{J_0 + \sum_{i=1}^{4} J_i \cos(i x_{1e} + \psi_i)}
	\end{bmatrix}= 
    \begin{bmatrix}
         0
        \\
        \frac{u_e - 1\cdot0-50\cdot\pi/3}{1.7+0.1\cdot cos(1\cdot\pi/3-0.04)+0.7\cdot cos(2\cdot \pi/3+2.9)+0.09\cdot cos(3\cdot \pi/3+2.8)+0.02\cdot cos(4\cdot \pi/3-2.6)}
    \end{bmatrix}
\end{align}
%

dal quale otteniamo
%
\begin{align}
	x_e := \begin{bmatrix}
	    \pi/3
        \\
        0
	\end{bmatrix},  \quad u_e = k\pi/3 \approx 52.36.\label{eq:equilibirum_pair}
\end{align}
%

Definiamo le variabili alle variazioni $\delta x$, $\delta u$ e $\delta y$ come 
%
\begin{align*}
	\delta x &= x-x_e, 
	\quad
	\delta u = u-u_e, 
	\quad
	\delta y = y-y_e.
\end{align*}
%

L'evoluzione del sistema espressa nelle variabili alle variazioni pu\`o essere approssimativamente 

descritta mediante il seguente sistema lineare
%
\begin{subequations}\label{eq:linearized_system}
\begin{align}
	\delta \dot{x} &= A\delta x + B\delta u
	\\
	\delta y &= C\delta x + D\delta u,
\end{align}
\end{subequations}
%

dove le matrici $A$, $B$, $C$ e $D$ vengono calcolate come
%
\begin{subequations}\label{eq:matrices}
\begin{align}
	A &= \begin{bmatrix}
\dfrac{\partial f_1(x,u)}{\partial x_1} & \dfrac{\partial f_1(x,u)}{\partial x_2} \\
\dfrac{\partial f_2(x,u)}{\partial x_1} & \dfrac{\partial f_2(x,u)}{\partial x_2}
\end{bmatrix}_{\substack{x = x_e \\ u = u_e}}
=
\begin{bmatrix}
0 & 1 \\
-\frac{kJ(x_{1e})}{[J(x_{1e})]^2} & -\frac{\beta}{J(x_{1e})}
\end{bmatrix} \approx \begin{bmatrix}
    0 & 1 
    \\
    -24.60 & -0.492
\end{bmatrix}
	\\
	B &= \begin{bmatrix}
\dfrac{\partial f_1(x,u)}{\partial u} \\
\dfrac{\partial f_2(x,u)}{\partial u}
\end{bmatrix}_{\substack{x = x_e \\ u = u_e}}
=
\begin{bmatrix}
0 \\
\frac{1}{J(x_{1e})}
\end{bmatrix} \approx
\begin{bmatrix}
    0
    \\
    0.492
\end{bmatrix}
	\\
	C &= \begin{bmatrix}
\dfrac{\partial h(x,u)}{\partial x_1} &
\dfrac{\partial h(x,u)}{\partial x_2}
\end{bmatrix}_{\substack{x = x_e \\ u = u_e}}
=
\begin{bmatrix}
1 & 0
\end{bmatrix}
	\\
	D &= \begin{bmatrix}
	    \dfrac{\partial h(x,u)}{\partial u} \end{bmatrix}_{\substack{x = x_e \\ u = u_e}}
= \begin{bmatrix}
    0
\end{bmatrix}.
\end{align}
\end{subequations}
%
\section{Calcolo Funzione di Trasferimento}

In questa sezione, andiamo a calcolare la funzione di trasferimento $G(s)$ dall'ingresso $\delta u$ all'uscita $\delta y$ 

mediante la seguente formula 
%
%
\begin{align}\label{eq:transfer_function}
G(s) = C(sI-A)^{-1} B+D \approx \frac{0.492}{s^{2} + 0.492\,s + 24.6}.
\end{align}
%

Dunque il sistema linearizzato~\eqref{eq:linearized_system} è caratterizzato dalla funzione di trasferimento~\eqref{eq:transfer_function} con due poli

complessi coniugati $p_{1,2} = -\sigma \pm j \omega_d= -0.2460 \pm 4.9536\,j$ e nessuno zero. 

Il sistema risulta asintoticamente stabile, in quanto la parte reale dei poli è negativa, pari a
\[
\Re\{p_{1,2}\} = -0.2460.
\]


Tuttavia, la presenza di poli complessi coniugati comporta una marcata componente oscillatoria nella 

risposta dinamica del sistema, che include un andamento sinusoidale smorzato.

Tale comportamento è ulteriormente accentuato dal basso valore del coefficiente di smorzamento 

relativo \(\xi\), definito come
\begin{align}
\xi = \frac{\sigma}{\sqrt{\sigma^2 + \omega_d^2}},
\end{align}


il quale, nel caso in esame, assume il valore
\begin{align}
\xi \approx \frac{0.2460}{\sqrt{0.2460^2 + 4.9536^2}} \approx 0.0496.
\end{align}

Il ridotto smorzamento relativo indica che il sistema è fortemente sottosmorzato; di conseguenza, 

la risposta al gradino è caratterizzata da una significativa sovraelongazione e da oscillazioni persistenti 

prima del raggiungimento del regime stazionario.

In Figura~\ref{fig:fig1} mostriamo
il corrispondente diagramma di Bode.

\begin{figure}[H]
    \centering
    \includegraphics[width=0.5\textwidth]{immagini/Figure_1.png}
    \caption{Diagramma di Bode della funzione di trasferimento $G(s)$.}
    \label{fig:fig1}
\end{figure} 

\label{sec:diagramma delle ampiezze}

Possiamo ricavare dal grafico i seguenti parametri caratteristici del sistema:
\begin{itemize}
    \item Pulsazione naturale: $\omega_n \approx 4.96 \ \text{rad/s}$;
    \item Guadagno statico: $K = 0.02 \; (\approx -34 \ \text{dB})$.
\end{itemize}

Tali valori possono essere verificati analiticamente a partire dai poli del sistema, fornendo una 

conferma della correttezza delle stime ottenute graficamente.

La pulsazione naturale è definita come
\begin{align}
\omega_n = \sqrt{\sigma^2 + \omega_d^2} \approx \sqrt{0.2460^2 + 4.9536^2} \approx 4.96 \ \text{rad/s},
\end{align}

in accordo con quanto osservato dal grafico.

Il guadagno statico \(K\) può invece essere ricavato come valore della funzione di trasferimento in 

continua, ovvero
\begin{align}
K = \lim_{s \to 0} G(s) = 0.02 \quad (\approx -34 \ \text{dB}),
\end{align}

coerentemente con l’ampiezza a bassa frequenza osservata nel diagramma di Bode.

Analizzando il diagramma di Bode, per quanto riguarda il diagramma delle ampiezze possiamo notare 

che:
\begin{itemize}
    \item A basse frequenze si ha un comportamento costante a $-34$ dB
    \item Intorno a $\omega_n$ si ha un picco di risonanza molto pronunciato dovuto al basso coefficiente di smorzamento. 
    \item Ad alte frequenze si ha una pendenza di $-40$ dB/decade 
\end{itemize}

Osserviamo invece che il diagramma delle fasi parte da $0^\circ$ nelle basse frequenze , per poi presentare 

una rapida variazione intorno a $\omega_n$ e arrivare infine asintoticamente a $-180^\circ$ nelle alte frequenze.


\section{Mappatura specifiche del regolatore}
\label{sec:specifications}
Le specifiche da soddisfare sono:
\begin{itemize}
	\item[1)] Errore a regime $|e_\infty| \le e^\ast = 0.04$ in risposta a un gradino $w(t) = W \cdot 1(t)$ e $d(t) = D \cdot 1(t)$ con ampiezze $W \le 1$ e $D \le 1$.
	\item[2)] Margine di fase $M_f \ge 30^\circ$.
	\item[3)] Sovraelongazione percentuale $S\% \le 10\%$.
    \item[4)] Tempo di assestamento $T_{a, 5} = 0.008\,\text{s}$.
    \item[5)] $d(t)$ attenuato di almeno  $30\,\text{dB}$ in $[0,\,0.5]$.
	\item[6)] $n(t)$ attenuato di almeno  $65\,\text{dB}$ in $[5\cdot 10^{4},\,5\cdot 10^{6}]$.
\end{itemize}
%
Andiamo ad effettuare la mappatura punto per punto delle specifiche richieste.
\begin{itemize}
	\item La specifica sull'errore a regime si traduce nella necessità di inserire un guadagno statico minimo.
	\item Le specifiche 2) e 3) sono entrambe sul margine di fase. Andiamo perciò a soddisfare la specifica più stringente, la 3), ottenendo così $M_f \ge 59.12^\circ$.
	\item Dalla specifica 4) sul tempo di assestamento $T_{a, 5}$ ricaviamo $\omega_c \geq \dfrac{300}{M_f\,T_{a,5}} \approx 634.35\ \mathrm{rad/s}$.
    \item Dalla specifica 5) ricaviamo $|L(j\omega)|_{dB} \geq 30$dB in $[0, 0.5]$.
\item Dalla specifica 6) ricaviamo $|L(j\omega)|_{dB} \leq -65$dB in $[5 \cdot 10^4, 5 \cdot10^6]$.
\end{itemize}
Pertanto, in Figura ~\ref{fig:fig2}, mostriamo il diagramma di Bode della funzione di trasferimento $G(s)$ con le zone proibite emerse dalla mappatura delle specifiche.


\begin{figure}[H]
    \centering
    \includegraphics[scale=0.5]{immagini/Figure_2.png}
    \caption{Diagramma di Bode di G(s) con zone proibite}
    \label{fig:fig2}
\end{figure}


Si osserva da Figura ~\ref{fig:fig2} che la funzione di trasferimento non soddisfa autonomamente le specifiche 

richieste. In particolare, il modulo presenta un andamento inadeguato nelle bande di 
frequenza 

critiche: a basse frequenze il guadagno risulta insufficiente per garantire l'attenuazione richiesta di 

$65\,\text{dB}$, mentre la fase del sistema evidenzia un comportamento incompatibile con il margine di fase 

$M_f \ge 59.12^\circ$ richiesto dalla specifica sulla sovraelongazione. Si rende pertanto indispensabile la 

progettazione di un regolatore opportuno.

\section{Sintesi del regolatore statico}
\label{sec:static_regulator}

In questa sezione progettiamo il regolatore statico $R_s(s)$ partendo dalle analisi svolte in sezione~\ref{sec:specifications}.


L'obiettivo primario di questa componente è soddisfare le specifiche di bassa frequenza: l'errore a 

regime e l'attenuazione del disturbo in uscita $d(t)$.

Per determinare il valore di $R_s$ è necessario garantire che il guadagno della funzione d'anello $L(s)$ sia 

sufficientemente elevato a basse frequenze, rispettando due vincoli distinti:
\begin{itemize}
    \item [1.]
    Per soddisfare la specifica 1) sull’errore a regime dobbiamo tenere conto della relazione $e_\infty = \frac{W +D}{1+ L(0)}$: un valore elevato di $L(0)$ permette infatti di ridurre l'errore a regime.
    Imponendo il vincolo $|e_\infty| \le e^\ast$ si trova che il guadagno di anello minimo necessario è il seguente 
    \begin {equation}
        \mu_{s,error} = \frac{D + W}{e^\ast} = 50
    \end{equation}

    \item[2.]
    Parallelamente, occorre rispettare la specifica 5) sull'attenuazione del disturbo di uscita $A_d$. Per ottenere l'attenuazione richiesta alla frequenza del disturbo $\omega_{d,max}$, il modulo della funzione di anello $|L(j \omega_{d,max})|$ deve superare il valore lineare 
     \begin {equation}
        mu_{s,dist} = 10^{\frac{A_d}{20}} \approx 31.62
    \end{equation}

\end{itemize}

Il valore di $R_s$ viene quindi determinato imponendo entrambi i vincoli e selezionando il più restrittivo 

tra i due:

\begin{subequations}
\begin{align}
  R_s = max(\frac{\mu_{s,error}}{|G(0)|}, \frac{\mu_{s,dist}}{|G(j\omega_{d,max}|}), \\
    \end{align}
\end{subequations}

Dunque, definiamo la funzione estesa $G_e(s) = R_s(s)G(s)$ e, in Figura \ref{fig:fig3}, mostriamo il suo diagramma 

di Bode per verificare se e quali zone proibite vengono ancora attraversate dal sistema.

\begin{figure}[H]
    \centering
    \includegraphics[scale=0.5]{immagini/Figure_3.png}
    \caption{Diagramma di Bode di $G_e(s)$ con zone proibite}
    \label{fig:fig3}
\end{figure}

Dalla Figura~\ref{fig:fig3} si osserva che l’introduzione del regolatore statico ha consentito di
soddisfare le 

specifiche a bassa frequenza. In particolare, la funzione estesa $G_e(s)$ rispetta ora
i vincoli sull’errore 

a regime e sull’attenuazione del disturbo di uscita $A_d$, come evidenziato
dal fatto che il modulo del 

diagramma di Bode si colloca al di fuori delle zone proibite nelle
bande di frequenza di interesse.

Tuttavia, la sola introduzione del regolatore statico non è sufficiente a garantire il rispetto
completo 

di tutte le specifiche di progetto. In particolare, $G_e(s)$ non soddisfa ancora i
vincoli imposti sul 

margine di fase $M_f$ e sul tempo di assestamento $T_{a,5}$.
\\

Si rende pertanto necessario introdurre un regolatore dinamico, in grado di modificare

opportunamente l’andamento della fase e del modulo della funzione di anello, così da fornire il

recupero di fase richiesto, aumentare la pulsazione di attraversamento e consentire il
raggiungimento 

delle prestazioni dinamiche e di robustezza richieste dalle specifiche.

\section{Sintesi del regolatore dinamico}
\label{sec:dynamic_regulator}

In questa sezione, progettiamo il regolatore dinamico $R_d(s)$. 
%
Dalle analisi fatte in Sezione~\ref{sec:static_regulator}, notiamo

di essere nello Scenario di tipo B (specifiche a bassa frequenza rispettate ma sistema instabile per 

margine di fase insufficiente e non abbastanza 
veloce ). Dunque, progettiamo $R_d(s)$ ricorrendo ad 

una rete anticipatrice e un filtro passa-basso per l'abbattimento del rumore nelle alte frequenze. 

Per determinare i parametri della rete utilizziamo le formule d’inversione, che permettono di imporre 

analiticamente il passaggio della funzione di anello per la pulsazione di attraversamento desiderata 

$\omega_c^\ast$ con il margine di fase specifico $M_f^\ast$. 
Tali valori sono stati scelti adottando un opportuno margine 

di cautela rispetto ai vincoli minimi di progetto, al fine di garantire una maggiore robustezza:
\begin{subequations}
\begin{align}
 M_f^\ast = M_{f,\text{spec}} + 10 \approx69,12\degree,
\qquad \omega_c^\ast = 1.1 \cdot \omega_{c,min} \approx697,7 ~\text{rad/s}.
\end{align}
\end{subequations}

Calcoliamo l'attenuazione $M^\ast$ e l'anticipo di fase $\varphi^\ast$ necessari a traslare la funzione d'anello $G_e(j\omega)$ 

verso la pulsazione di attraversamento $\omega_c^\ast$ e il margine $M_f^\ast$ prefissati:
\begin{equation}
M^\ast = \frac{1}{|G_e(j\omega_c^\ast)|}, \quad \varphi^\ast = M_f^\ast - 180^\circ - arg( G_e(j\omega_c^\ast))
\end{equation}


A questo punto calcoliamo i parametri della rete $\tau$ e $\alpha$ con le formule di inversione:
\begin{equation}
\tau = \frac{M^\ast - \cos(\varphi^\ast)}{\omega_c^\ast \sin(\varphi^\ast)}, \quad \alpha\tau = \frac{\cos(\varphi^\ast) - \frac{1}{M^\ast}}{\omega_c^\ast \sin(\varphi^\ast)}
\end{equation}

Otteniamo così il modulo richiesto:
\begin{equation}
    R_d(s) = \frac{1 + \tau s}{1 + \alpha \tau s}.
\end{equation}

Infine per rispettare il vincolo sul disturbo $n(t)$, aggiungiamo un polo 
ad alta frequenza:
\[
R_{\text{hf}}(s) = \frac{1}{1 + \tfrac{s}{9\cdot 10^4}},
\]

La forma finale del regolatore dinamico diventa perciò:
\begin{equation}
    R_d(s) = 
    \frac{1 + \tau s}{1 + \alpha\tau s}
    \cdot
    \frac{1}{1 + \tfrac{s}{9\cdot 10^4}}.
\end{equation}

La funzione di anello chiuso risulta:
\[
L(s) = R_d(s)\,G_e(s).
\]



In Figura \ref{fig:fig4}, mostriamo il diagramma di Bode della funzione d'anello $L(s)$

\begin{figure}[H]
    \centering
    \includegraphics[scale=0.5]{immagini/Figure_4.png}
    \caption{Diagramma di Bode della funzione di anello $L(s) = R_d(s)G_e(s)$.}
    \label{fig:fig4}
\end{figure}

Dalla Figura ~\ref{fig:fig4} si osserva che la funzione di anello $L(s)$ evita correttamente tutte le zone
proibite 

individuate in fase di mappatura delle specifiche. In particolare, il modulo a bassa frequenza
rimane 

sufficientemente elevato da garantire il rispetto delle specifiche sull’errore a regime e
sull’attenuazione 

del disturbo $d(t)$, mentre la pulsazione di attraversamento risulta compatibile con
il vincolo sul tempo 

di assestamento $T_{a,5}=0.008$ s.


Grazie all’azione della rete anticipatrice, la fase viene opportunamente incrementata in corrispondenza

della nuova pulsazione $\omega_c^\ast$, consentendo di ottenere un margine di fase
coerente con le specifiche di 

stabilità relativa e tale da assicurare una risposta sufficientemente smorzata, con sovraelongazione 

inferiore al $10\%$.

Infine, l’introduzione del polo ad alta frequenza determina un rapido decadimento del modulo alle 

alte frequenze, permettendo di soddisfare il requisito di attenuazione del rumore di misura $n(t)$ di 

almeno $65$ dB nella banda assegnata, senza compromettere la stabilità del sistema.

\section{Test sul sistema linearizzato}

In questa sezione testiamo l'efficacia del controllore progettato sul sistema linearizzato. 
In primo 

luogo analizziamo le risposte del sistema al riferimento a gradino $w(t)$ e ai disturbi $n(t)$ e $d(t)$, definiti 

come segue:
\[
w(t) = \mathbf{1}(t),
\qquad
d(t) = \sum_{k=1}^{4} \sin(0.1\,k\,t),
\qquad
n(t) = \sum_{k=1}^{4} \sin\!\left(5 \cdot 10^{4}\, k\, t\right).
\]

Successivamente, consideriamo la risposta del sistema linearizzato in termini dell’uscita complessiva 

$y_{\text{tot}}$, ottenuta come somma delle singole risposte, facendo riferimento al principio di sovrapposizione 

degli effetti.

\vspace{1em} 
Innanzitutto mostriamo in Figura~\ref{fig:fig5} lo schema a blocchi del sistema in anello chiuso.
\begin{figure}[H]
    \centering
    \includegraphics[scale=0.5]{immagini/Figure_5.png}
    \caption{Schema a blocchi del sistema lineare}
    \label{fig:fig5}
\end{figure}

\vspace{1em} 

Di seguito in Figura~\ref{fig:fig6} è riportata la risposta del sistema linearizzato in anello chiuso a fronte di un 

ingresso a gradino unitario $w(t) = 1$, ottenuta simulando $y_w(t) = F(s) \cdot w(t)$.

\begin{figure}[H]
    \centering
    \includegraphics[scale=0.5]{immagini/Figure_6.png}
    \caption{Risposta del sistema al riferimento a gradino unitario}
    \label{fig:fig6}
\end{figure}

Si nota che il sistema raggiunge il valore di regime unitario con un tempo di assestamento inferiore a 

$T_{a,5} = 0.008$ s, soddisfacendo la specifica 4). La risposta presenta una sovraelongazione contenuta, al 

di sotto del limite del $10\%$ imposto dalle specifiche. L'errore a regime è praticamente nullo, perciò 

anche la specifica 1) $e_{reg} < 0.04$ viene rispettata.

\vspace{1em} 

Inoltre possiamo notare dalla seguente Figura~\ref{fig:fig7}  che i disturbi in attuazione $d(t)$, costituiti da una 

somma di quattro armoniche a bassa frequenza ($\omega_d = 0.1$ rad/s e multipli), vengono efficacemente 

attenuati dalla funzione di sensitività $S(s)$, come mostrato nella risposta $y_d(t) = S(s) \cdot d(t)$.

\begin{figure}[H]
    \centering
    \includegraphics[scale=0.5]{immagini/Figure_7.png}
    \caption{Risposta del sistema al disturbo in attenuazione.}
    \label{fig:fig7}
\end{figure}
Osserviamo che l'attenuazione ottenuta è superiore ai $30$ dB richiesti dalla specifica 5) nella banda 

$[0, 0.5]$ rad/s.

\vspace{1em} 

In seguito, valutiamo la capacità del sistema di attenuare il rumore di misura $n(t)$, composto da 

quattro armoniche ad alta frequenza ($\omega_n = 5 \cdot 10^4$ rad/s e multipli). La risposta $y_n(t) = -F(s) \cdot n(t)$ 

è mostrata in Figura~\ref{fig:fig8}.

\begin{figure}[H]
    \centering
    \includegraphics[scale=0.5]{immagini/Figure_8.png}
    \caption{Risposta del sistema al rumore di misura}
    \label{fig:fig8}
\end{figure}

Notiamo che l'uscita $y_n(t)$ presenta un'attenuazione superiore ai $65$ dB richiesti dalla specifica 6) 

nella banda $[5 \cdot 10^4, 5 \cdot 10^6]$ rad/s. Questo è dovuto al polo in alta frequenza inserito nel regolatore, 

che si comporta con un filtro passa-basso che elimina tutte le frequenze superiori alla banda passante

del sistema, in questo modo il rumore non si propaga in uscita.

\vspace{1em} 

Per valutare le prestazioni complessive del sistema in condizioni operative realistiche, consideriamo 

la risposta totale ottenuta come somma dei tre contributi:
\begin{equation}
    y_{tot}(t) = y_w(t) + y_d(t) + y_n(t)
\end{equation}

\begin{figure}[H]
    \centering
    \includegraphics[scale=0.5]{immagini/Figure_9.png}
    \caption{Risposta totale del sistema in presenza di riferimento, disturbo e rumore}
    \label{fig:fig9}
\end{figure}

Da Figura \ref{fig:fig9} notiamo che l'uscita totale soddisfa i vincoli relativi alla sovraelongazione percentuale e 

al tempo di assestamento. È importante notare che, nonostante la presenza di disturbi a bassa 

frequenza e rumore ad alta frequenza, l'effetto combinato sulla risposta totale risulta trascurabile 

grazie alle elevate attenuazioni garantite rispettivamente dalla 

funzione di sensitività $S(s)$ e dalla funzione di sensitività complementare $F(s)$ nelle rispettive bande 

di frequenza critiche.

\section{Test sul sistema non lineare}

In questa sezione, testiamo l’efficacia del controllore progettato sul modello non lineare con l’idea di 

applicare al sistema reale il regolatore sintetizzato sul modello linearizzato. Il controllore $R(s)$ opera 

sulle variabili alle variazioni: riceve in ingresso l’errore
\[
e(t) = \tilde y(t) - w(t) \approx \Delta y(t) - w(t)
\]

e genera come uscita una variazione di controllo
\[
\tilde u(t) \approx \Delta u(t).
\]

Poiché il sistema non lineare lavora con le variabili originali, l’ingresso effettivamente applicato è 

ricostruito come
\[
u(t) = u_e + \tilde u(t).
\]

In questo modo il regolatore risulta coerente con il sistema reale, purché il funzionamento avvenga in 

un intorno del punto di equilibrio, dove l’approssimazione di linearizzazione è valida. \\

In Figura \ref{fig:fig10}, mostriamo lo schema a blocchi del sistema non lineare in anello chiuso.

\begin{figure}[H]
    \centering
    \includegraphics[scale=0.5]{immagini/Figure_10.png}
    \caption{Schema a blocchi del sistema in anello chiuso.}
    \label{fig:fig10}
\end{figure}


Di seguito è riportato, in Figura \ref{fig:fig11}, l'andamento temporale dell'uscita del sistema in risposta a un 

ingresso a gradino di ampiezza unitaria considerando i disturbi $n(t)$ e $d(t)$.

\begin{figure}[H]
    \centering
    \includegraphics[scale=0.5]{immagini/Figure_11.png}
    \caption{\centering Risposta del sistema non lineare in anello chiuso: andamento dell'uscita in presenza del riferimento e dei disturbi.}
    \label{fig:fig11}
\end{figure}


Si nota che l'uscita del sistema raggiunge il valore di regime in modo compatibile con le specifiche di 

progetto sull'errore a regime poste nella fase di sintesi del controllore.

La coerenza riscontrata tra i risultati delle simulazioni sul sistema linearizzato e quelle sul sistema 

non linearizzato conferma l'efficacia della linearizzazione nell'intorno del punto di equilibrio $\theta_e = \pi/3$. 

Il tempo di assestamento rimane rapido, il sistema non presenta sovraelongazioni marcate e il rumore 

ad alta frequenza $n(t)$ risulta ben attenuato con un impatto sull'uscita finale trascurabile.
La stabilità 

globale del sistema è quindi preservata anche in presenza della variazione periodica del momento 

d'inerzia $J(\theta)$.\\

Come da richiesta, fissato il riferimento a gradino
$w(t)=1(t)$, è stato analizzato il comportamento 

del sistema non lineare in anello chiuso al variare delle condizioni iniziali dello stato nell’intorno del 

punto di equilibrio
\[
x_e=\begin{bmatrix}\theta_e\\0\end{bmatrix}, \qquad \theta_e=\pi/3.
\]

Il valore di regime dell’uscita in presenza del gradino unitario risulta, nel caso base \(x(0)=x_e\), pari a
\[
y_{\infty,\text{base}} \simeq 2.043\ \text{rad},
\]

con un errore a regime \(|e_\infty|\approx 0.004\ \text{rad}\), inferiore alla soglia di tolleranza assegnata di \(0.04\ \text{rad}\).

Sono quindi state esplorate le condizioni iniziali dello stato
\[
x(0)=
\begin{bmatrix}
\theta(0)\\
\omega(0)
\end{bmatrix}
=
\begin{bmatrix}
\theta_e+\Delta\theta\\
\omega_e +\Delta\omega
\end{bmatrix},
\]

verificando, tramite simulazioni del modello non lineare in anello chiuso, la convergenza dell’uscita 

allo stesso valore di regime e il rispetto della specifica sull’errore a regime.

In particolare, mantenendo \(\omega(0)=0\), e tenendo in considerazione che $\theta_e = \pi/3$ si è osservata 

convergenza e rispetto della condizione \(|e_\infty|<0.04\ \text{rad}\) per:
\[
\theta(0)\in[\theta_e-29,\ \theta_e+28]\ \text{rad}.
\]


Analogamente, mantenendo \(\theta(0)=\theta_e\), e tenendo in considerazione che $\omega_e = 0$, si è verificato che il 

sistema converge correttamente per:
\[
\omega(0)\in[\omega_e-30,\ \omega_e+ 30]\ \text{rad/s}.
\]

Tali risultati forniscono una stima numerica della regione di attrazione locale del punto di equilibrio 

in anello chiuso nel piano delle condizioni iniziali \((\theta(0),\omega(0))\), all’interno della quale l’uscita del sistema 

non lineare converge al valore di regime associato al riferimento applicato e l’errore a regime soddisfa 

la specifica assegnata.

Come da richiesta, supponendo un riferimento $w(t) = W \cdot1(t)$, è stato esplorato il range di ampiezza 

di
riferimenti a gradino tali per cui il controllore rimane efficace sul sistema non lineare. Si è appurato 

che l'ampiezza debba rientrare nel seguente intervallo $ W \in [ -3 ; 3]$ rad affinché l'errore a regime $|e_{\infty}|$ 

rimanga al di sotto della soglia di tolleranza pari a 0.04 rad.


Concludendo, possiamo affermare che il regolatore progettato si è dimostrato robusto, garantendo 

il soddisfacimento delle specifiche non solo nel caso del modello linearizzato, ma anche a fronte della 

dinamica non lineare completa del sistema reale.
\\

\section{Conclusioni}

Il progetto ha affrontato il problema del controllo di un meccanismo non lineare attuato. A partire 

dal modello non lineare, il sistema è stato espresso in forma di stato e linearizzato intorno a un punto 

di equilibrio assegnato, ottenendo così un sistema lineare tempo invariante per la sintesi del controllore.

L’analisi del sistema linearizzato ha evidenziato la necessità dell’impiego di un regolatore in 

retroazione. La progettazione è stata condotta in frequenza, traducendo le specifiche temporali e di 

reiezione dei disturbi in vincoli sul diagramma di Bode della funzione di anello.

Il controllore è stato progettato combinando un regolatore statico, volto a soddisfare le specifiche 

sull’errore a regime e sull’attenuazione dei disturbi a bassa frequenza, con un regolatore dinamico 

anticipatore e un polo ad alta frequenza, al fine di garantire un adeguato margine di fase, un tempo 

di assestamento compatibile e una forte attenuazione del rumore di misura.

Le simulazioni sul sistema linearizzato confermano il rispetto delle specifiche assegnate, mostrando 

una buona rapidità di risposta e un’efficace reiezione delle perturbazioni. L’applicazione del controllore 

al modello non lineare evidenzia che, nonostante la natura locale della linearizzazione, i risultati hanno 

confermato una marcata coerenza tra il comportamento del modello semplificato e quello della 

dinamica reale. I test condotti hanno inoltre dimostrato la robustezza del regolatore a fronte di 

variazioni del riferimento, evidenziando la capacità del controllo di gestire efficacemente le non 

linearità intrinseche del meccanismo in un ampio intervallo operativo ($W \in [-3, 3]$ rad).

\end{document}
